\documentclass[a4paper,12pt]{article}
\usepackage{geometry}
\geometry{margin=1in}
\usepackage{amsmath}
\usepackage{graphicx}
\usepackage{enumitem}
\usepackage{fancyhdr}
\usepackage{titlesec}
\usepackage{setspace}

% Set up header and footer
\pagestyle{fancy}
\fancyhf{}
\fancyhead[L]{CBS Examination, Dec 2018}
\fancyhead[R]{Computer Networks (CE-204)}
\fancyfoot[C]{\thepage}

% Title formatting
\titleformat{\section}[block]{\large\bfseries}{\thesection.}{1em}{}
\titleformat{\subsection}[block]{\normalsize\bfseries}{\thesubsection.}{1em}{}

% Set spacing
\setstretch{1.2}

\title{\vspace{-2cm}CBS Examination, Dec 2018\\\textbf{Computer Networks (CE-204)}}
\author{}
\date{}

\begin{document}

\maketitle

\noindent \textbf{Time: 3 Hours} \hfill \textbf{Max. Marks: 60}

\section*{Instructions:}
\begin{enumerate}[label=\arabic*.]
    \item It is compulsory to answer all the questions (2 marks each) of Part-A in short.
    \item Answer any four questions from Part-B in detail.
    \item Different sub-parts of a question are to be attempted adjacent to each other.
\end{enumerate}

\section*{Part-A}

\begin{enumerate}[label=(\alph*)]
    \item List two characteristics of a broadband network. \hfill (2)
    \item What is HTTP? \hfill (2)
    \item Mention the advantages of wireless LAN. \hfill (2)
    \item What is the difference between network layer delivery and transport layer delivery? \hfill (2)
    \item Name some parameters for determining the quality of service at the transport layer. \hfill (2)
    \item Discuss basic functions of the data link layer in brief. \hfill (2)
    \item What is the effect of packet size on transmission time? \hfill (2)
    \item What is dotted decimal notation? \hfill (2)
    \item What is the difference between a port address, a logical address, and a physical address? \hfill (2)
    \item Differentiate between a point-to-point direct link and a multipoint guided configuration. \hfill (2)
\end{enumerate}

\section*{Part-B}
\textbf{Q2:} 
\begin{enumerate}[label=(\alph*)]
    \item What is a LAN? Can a LAN have routers? What is the difference between a packet and a frame? \hfill (5)
    \item Create a diagram of the NRZ, NRZI, and Manchester encodings for the bit pattern 11010011. \hfill (5)
\end{enumerate}

\textbf{Q3:}
\begin{enumerate}[label=(\alph*)]
    \item Describe sliding window flow control. \hfill (5)
    \item Explain CRC coding mechanism considering the following message and generator polynomial. \hfill (5)
    \begin{itemize}
        \item Message: $M=110010101100101010$
        \item Generator Polynomial: $G=1010$
    \end{itemize}
    Show the working of CRC on this data.
\end{enumerate}

\textbf{Q4:}
\begin{quote}
    Explain the advantages of IPv6 over IPv4. A large number of consecutive IP addresses are available starting at 198.16.0.0. Suppose that four organizations A, B, C, and D request 4000, 2000, 4000, and 8000 addresses, respectively, in that order.  
    \begin{enumerate}[label=(\roman*)]
        \item First IP address assigned
        \item Last IP address assigned
        \item Mask in the W.X.Y.Z / S notation
    \end{enumerate}
\end{quote}

\textbf{Q5:}
\begin{enumerate}[label=(\alph*)]
    \item Compare the TCP header and UDP header. List the fields in the TCP header that are missing from the UDP header. Give the reason for their absence. \hfill (5)
    \item What is the Address Resolution Protocol? How is mapping performed between an IP address and a MAC address? \hfill (5)
\end{enumerate}

\textbf{Q6:}
\begin{enumerate}[label=(\alph*)]
    \item What is the difference between open-loop congestion control and closed-loop congestion control? \hfill (5)
    \item Explain the Distance Vector Routing Protocol? What is its major drawback? \hfill (5)
\end{enumerate}

\textbf{Q7:}
Differentiate between the following: \hfill (10)
\begin{enumerate}[label=(\alph*)]
    \item Pure ALOHA and Slotted ALOHA
    \item Packet Switched Networks and Circuit Switched Networks
\end{enumerate}

\end{document}
